\documentclass[11pt]{article}
\usepackage[a4paper,margin=27mm]{geometry}
\usepackage{amsmath,amssymb,mathtools}
\usepackage{fontspec,unicode-math}
\setmainfont{Latin Modern Roman}
\setmathfont{Latin Modern Math}
\usepackage{xurl}
\usepackage[hidelinks]{hyperref}
\newcommand{\sref}[2]{\hyperlink{source:#1:#2}{[S#2]}}
\newcommand{\cataloguescope}[1]{\paragraph{Mathematical question.}#1}
\hypersetup{pdftitle={Problem statement}}
\begin{document}
\section{Are there infinitely many Fibonacci primes}\label{statement}

\subsection{Definitions and mathematical statement}

\paragraph{Shared notation.}$\mathbb N=\{1,2,\ldots\}$, and $\mathbb Z,\mathbb Q,\mathbb R,\mathbb C$ denote the integers, rationals, reals and complex numbers. Any other conventions are specified locally.


Define Fibonacci numbers by $F_0=0$, $F_1=1$ and $F_{n+2}=F_{n+1}+F_n$ for $n\ge0$. A Fibonacci prime is a prime value $F_n$.
\paragraph{Statement recorded in the supplied source.}
Are there infinitely many integers $n\ge0$ for which $F_n$ is prime? \sref{statement}{2}
\subsection{Short English statement}
Do infinitely many Fibonacci numbers take prime values?
\subsection{Sources}
\begin{description}
\item[\hypertarget{source:statement:1}{S1}] ULAM.AI and the UnsolvedMath Contributors, UnsolvedMath: A Curated Collection of Open Mathematics Problems, catalogue entry. CC BY 4.0. \url{https://huggingface.co/datasets/ulamai/UnsolvedMath}.
\item[\hypertarget{source:statement:2}{S2}] Wikipedia contributors, \emph{List of unsolved problems in mathematics}, revision 1366636767 (2026) Number theory, Are there infinitely many Fibonacci primes. \url{https://en.wikipedia.org/w/index.php?title=List_of_unsolved_problems_in_mathematics&oldid=1366636767#Number_theory}.
\item[\hypertarget{source:statement:3}{S3}] OEIS Foundation and contributors, \emph{The On-Line Encyclopedia of Integer Sequences}, A000045: Fibonacci numbers: F(n) = F(n-1) + F(n-2) with F(0) = 0 and F(1) = 1., official data snapshot \texttt{693b5499c9d2} (2026) Name, definition and comments. \url{https://raw.githubusercontent.com/oeis/oeisdata/693b5499c9d248f570dd2d979e8bd1d7eecf689b/seq/A000/A000045.seq}.
\item[\hypertarget{source:statement:4}{S4}] OEIS Foundation and contributors, \emph{The On-Line Encyclopedia of Integer Sequences}, A001605: Indices of prime Fibonacci numbers., official data snapshot \texttt{693b5499c9d2} (2026) Name, definition and comments. \url{https://raw.githubusercontent.com/oeis/oeisdata/693b5499c9d248f570dd2d979e8bd1d7eecf689b/seq/A001/A001605.seq}.
\end{description}
\label{end:statement}
\end{document}
